\section{Introduction}
\label{sec:introduction}

Electroencephalography (EEG) records neural activity at the scalp with millisecond resolution, but the same electrodes also capture potentials generated by blinks and eye movements. These ocular artifacts are often larger than the EEG activity of interest and spread across many channels. Rejecting every affected segment can remove a substantial fraction of a recording, while aggressive correction can attenuate event-related potentials, oscillatory activity, or spatial structure. Ocular-artifact removal is therefore a restoration problem: the method must reduce the contribution of the eyes while preserving the neural signal that overlaps it in time and frequency.

Existing approaches address this problem through EOG regression, adaptive filtering, blind source separation, artifact-subspace methods, and supervised neural networks \cite{gratton1983ocular,semlitsch1986ocular,jung2000bss,chang2020asr,zhang2021eegdenoisenet,yu2022deepseparator}. Classical regression is attractive when electrooculography (EOG) is recorded because its correction coefficients have a direct interpretation, but fixed coefficients do not capture changes in electrode placement and ocular propagation. Learned denoisers provide more flexible mappings, yet a single population model must average over these participant- and session-dependent differences. Diffusion models offer a conditional generative formulation for EEG restoration \cite{huang2025eegdfus}, but conditioning on the corrupted waveform alone does not identify how a particular participant's eye activity reaches the scalp electrodes.

The key observation behind this work is that the participant-dependent quantity needed for ocular correction is not an identity label. It is the EOG-to-EEG propagation measured during the recording session. Calibration-based ocular correction has long estimated this propagation from simultaneous EEG and EOG \cite{gratton1983ocular,kobler2020sgeyesub}; we use the same physical cue to adapt a population conditional diffusion model. A short calibration segment is kept separate from evaluation data and used to estimate a propagation matrix. At evaluation time, the recorded EOG and this matrix produce an ocular-artifact estimate that guides the denoiser. Empirical Bayes shrinkage stabilizes the estimate and returns low-reliability calibration to the population estimate rather than forcing an individual correction.

This design yields a single subject-calibrated EOG-guided diffusion method. The neural-signal model is shared across participants, whereas the ocular observation model is calibrated within each participant--session--task cell. The separation permits direct controls: population calibration tests whether an individual estimate is useful, calibration from another participant tests specificity, and a zero-anchor control measures the contribution of the EOG-derived waveform guide while retaining the same network and compact calibration state. Repeated diffusion samples are also used to form empirical predictive intervals, with part of their width assigned to uncertainty in the estimated propagation matrix.

We evaluate the method on a development cohort and a held-out participant cohort, using paired semi-simulation for waveform error and natural recordings for EOG attenuation and signal-preservation measures. Matched calibration improves over both population calibration and the zero-anchor control in development, and the zero-anchor comparison is repeated without method selection in held-out participants. A controlled injected sensitivity analysis examines whether a support-fitted representation can carry the same participant-matched calibration information through a different interface. Calibration-aware predictive intervals improve empirical coverage without increasing the continuous ranked probability score relative to the unadjusted sample distribution. The experiments support a practical separation of roles: a population-trained diffusion model represents common EEG structure, while a short EEG--EOG calibration adapts the artifact pathway.
